% Image credits for the photographs and AI illustrations of Book 4
% (University Chemistry, Year 3). Machine-readable license records live in
% images/book4/CREDITS.md; keep the two files in sync. The Book 4 writer
% owns this file: add one \item per photograph (in an itemize, after the
% first paragraph) as photographs land.
\chapter*{Kredit Gambar}

Gambar-gambar dalam buku ini adalah karya asli. Foto-foto, jika dipakai,
dipakai dengan lisensi bebasnya masing-masing, dengan terima kasih kepada
para pembuatnya; tautan sumber dan URL lisensi lengkap untuk setiap foto
tercantum dalam \texttt{images/book4/CREDITS.md} di repositori buku ini.
Piktogram bahaya adalah piktogram Sistem Harmonisasi Global untuk
Klasifikasi dan Pelabelan Bahan Kimia, sebagaimana digambar untuk paket
\texttt{ghsystem} dari \TeX\ Live.

\bigskip
Di samping foto-foto dan gambar-gambar asli, beberapa gambar adalah
ilustrasi buatan AI, yang dihasilkan dengan pembangkit gambar OpenAI dari
perintah yang ditulis oleh para penyunting buku ini, dan diperiksa
ketepatan kimianya sebelum dimuat. Perintah lengkap untuk setiap gambar
yang dihasilkan dicatat dalam \texttt{images/book4/ai/PROMPTS.md} di
repositori.
\bigskip
\noindent Foto-foto, menurut bab:
\begin{itemize}\small
  \item Bab~1: Erwin Schr\"odinger, 1933, Yayasan Nobel, domain publik
        (Wikimedia Commons).
  \item Bab~4: kristal salju, Wilson A. Bentley (sekitar 1902), domain publik
        (Wikimedia Commons).
  \item Bab~6: antena radio di sebuah dataran tinggi, ESO/B.~Tafreshi
        (twanight.org), CC BY 4.0 (Wikimedia Commons).
  \item Bab~7: fluorit dalam cahaya siang dan di bawah sinar ultraviolet, oleh Masha
        Milshina, CC BY 4.0 (Wikimedia Commons).
  \item Bab~8: sebuah laboratorium NMR, oleh Shandchem, CC BY 2.0 (Wikimedia
        Commons).
  \item Bab~9: William Lawrence Bragg, Yayasan Nobel, domain publik
        (Wikimedia Commons).
  \item Bab~10: makam Ludwig Boltzmann, oleh Feldkurat Katz,
        CC BY-SA 4.0 (Wikimedia Commons).
  \item Bab~13: spiral Belousov--Zhabotinsky, oleh Stephen Morris, CC BY 2.0
        (Wikimedia Commons).
  \item Bab~14: lubang ozon Antartika, NASA, domain publik (Wikimedia
        Commons).
  \item Bab~16: mikrograf elektron sarang lebah sebuah konverter katalitik, oleh
        Galadrid, CC BY-SA 4.0 (Wikimedia Commons).
  \item Bab~17: berkas sinar matahari dan efek Tyndall, oleh Kevin c higgins, CC BY-SA 4.0
        (Wikimedia Commons).
  \item Bab~18: kristal rubi dan zamrud, oleh Robert M. Lavinsky, CC BY-SA 3.0
        (Wikimedia Commons).
  \item Bab~20: kristal ferosena, oleh Andrea Sella, CC BY 4.0 (Wikimedia
        Commons).
  \item Bab~22: bola kristal tunggal silikon, oleh Sebastian Wallroth, domain
        publik (Wikimedia Commons).
  \item Bab~23: sekuntum bunga pada aerogel silika, NASA/JPL, domain publik; piala
        Lycurgus, oleh Marie-Lan Nguyen, CC BY 2.5 (Wikimedia Commons).
  \item Bab~24: seekor belangkas Atlantik, oleh Kaldari, CC0 (Wikimedia
        Commons).
  \item Bab~27: bunga piretrum, oleh Soramimi, CC BY-SA 4.0 (Wikimedia
        Commons).
  \item Bab~30: kulit kayu pohon yew Pasifik, oleh JOE BLOWE (theforestprimeval),
        CC BY-SA 2.0 (Wikimedia Commons).
  \item Bab~32: ledakan alga yang dilihat oleh satelit Landsat 8, USGS dan NASA,
        domain publik.
  \item Bab~33: manifold ganda vakum dan gas inert, oleh Polimerek, CC BY-SA
        3.0 (Wikimedia Commons).
\end{itemize}
\clearpage
